% GENERATED FILE. Edit canonical lessons, then run npm run generate:books.
% canonical-source-hash: fnv1a64:be63765b8ef6bd52
% canonical-lessons: RU-C68-day, RU-C68-week, RU-C68-month, RU-C68-year, RU-C68-evening

\chapter{Day, Week, Month, Year, Evening}
\label{ch:ru-day-week-month-year-evening}
\hlchaptermodality{68}

\begin{chapteropening}
\textbf{By the end of this chapter:} \emph{I can name a day, a week, a month, a year and the evening.}

Complete the last lesson of chapter 68: I can name a day, a week, a month, a year and the evening.
\end{chapteropening}

\section[\texorpdfstring{\ru{день} (den)}{день (den)}]{\ru{день} (den) — a day}

\label{lesson:RU-C68-day}



\begin{quote}
Before the new one: say the Russian for cheap, then the Russian for badly.
\end{quote}

\subsection*{You'll want to know: \ru{день}}
\textbf{\ru{день}} (\emph{den}) — \textquotedblleft{}a day\textquotedblright{}.

\begin{cousinweb}
\textbf{\ru{день}} $\leftarrow$ Proto-Slavic \emph{*dĭnĭ} — inherited. Vasmer lists Sanskrit \emph{dínam}, Lithuanian \emph{dienà} and Latin \emph{diēs}, all \textquotedblleft{}day\textquotedblright{}, from the PIE root \emph{*dyew-}, \textquotedblleft{}to shine\textquotedblright{}.
\end{cousinweb}

\begin{grammarlens}[title={What you've built}]
One more word for this chapter.
\end{grammarlens}

\subsection*{Your turn}
\begin{itemize}
\raggedright
  \item \emph{Say it:} \emph{den}
  \item \emph{Say it:} \emph{den}, once more
  \item \emph{Say it:} say \emph{den}
  \item \emph{Recall:} say the Russian for cheap, then the Russian for badly, then say \emph{den} again
\end{itemize}

\begin{tcolorbox}[breakable,colback=teal!4,colframe=teal!35!black,title={Before you move on}]
What does \ru{день} mean? (\textquotedblleft{}A day\textquotedblright{}.) Say it once more.
\end{tcolorbox}
\section[\texorpdfstring{\ru{неделя} (nedélya)}{неделя (nedélya)}]{\ru{неделя} (nedélya) — a week}

\label{lesson:RU-C68-week}



\begin{quote}
Before the new one: say the Russian for badly, then the Russian for a day.
\end{quote}

\subsection*{You'll want to know: \ru{неделя}}
\textbf{\ru{неделя}} (\emph{nedélya}) — \textquotedblleft{}a week\textquotedblright{}.

\begin{cousinweb}
\textbf{\ru{неделя}} is \emph{\ru{не}}, \textquotedblleft{}not\textquotedblright{}, with \emph{\ru{дело}}, \textquotedblleft{}work\textquotedblright{}: first \textquotedblleft{}the day of no work\textquotedblright{}, Sunday, and then the week that runs from it (Vasmer). Vasmer thinks it most likely copies Latin \emph{fēria}, \textquotedblleft{}holiday\textquotedblright{}.
\end{cousinweb}

\begin{grammarlens}[title={What you've built}]
One more word for this chapter.
\end{grammarlens}

\subsection*{Your turn}
\begin{itemize}
\raggedright
  \item \emph{Say it:} \emph{nedélya}
  \item \emph{Say it:} \emph{nedélya}, once more
  \item \emph{Say it:} say \emph{nedélya}
  \item \emph{Recall:} say the Russian for badly, then the Russian for a day, then say \emph{nedélya} again
\end{itemize}

\begin{tcolorbox}[breakable,colback=teal!4,colframe=teal!35!black,title={Before you move on}]
What does \ru{неделя} mean? (\textquotedblleft{}A week\textquotedblright{}.) Say it once more.
\end{tcolorbox}
\section[\texorpdfstring{\ru{месяц} (mésyats)}{месяц (mésyats)}]{\ru{месяц} (mésyats) — a month}

\label{lesson:RU-C68-month}



\begin{quote}
Before the new one: say the Russian for a day, then the Russian for a week.
\end{quote}

\subsection*{You'll want to know: \ru{месяц}}
\textbf{\ru{месяц}} (\emph{mésyats}) — \textquotedblleft{}a month\textquotedblright{}.

\begin{cousinweb}
\textbf{\ru{месяц}} $\leftarrow$ Proto-Slavic \emph{*měsęcĭ}, \textquotedblleft{}moon; month\textquotedblright{} — inherited. Vasmer lists Sanskrit \emph{mās-}, Greek \emph{m\'{\={e}}n} and Latin \emph{mēnsis}, \textquotedblleft{}month\textquotedblright{}, and Lithuanian \emph{mėnuo}, \textquotedblleft{}moon\textquotedblright{}: English \emph{moon} and \emph{month} are cousins. A month is one round of the moon.
\end{cousinweb}

\begin{grammarlens}[title={What you've built}]
One more word for this chapter.
\end{grammarlens}

\subsection*{Your turn}
\begin{itemize}
\raggedright
  \item \emph{Say it:} \emph{mésyats}
  \item \emph{Say it:} \emph{mésyats}, once more
  \item \emph{Say it:} say \emph{mésyats}
  \item \emph{Recall:} say the Russian for a day, then the Russian for a week, then say \emph{mésyats} again
\end{itemize}

\begin{tcolorbox}[breakable,colback=teal!4,colframe=teal!35!black,title={Before you move on}]
What does \ru{месяц} mean? (\textquotedblleft{}A month\textquotedblright{}.) Say it once more.
\end{tcolorbox}
\section[\texorpdfstring{\ru{год} (god)}{год (god)}]{\ru{год} (god) — a year}

\label{lesson:RU-C68-year}



\begin{quote}
Before the new one: say the Russian for a week, then the Russian for a month.
\end{quote}

\subsection*{You'll want to know: \ru{год}}
\textbf{\ru{год}} (\emph{god}) — \textquotedblleft{}a year\textquotedblright{}.

\begin{cousinweb}
\textbf{\ru{год}} $\leftarrow$ Proto-Slavic \emph{*godŭ}, \textquotedblleft{}a fitting time, a season\textquotedblright{} — inherited; Old Russian \emph{godŭ} meant \textquotedblleft{}time, term\textquotedblright{} (Vasmer). Vasmer relates it to Latvian \emph{gadīgs}, \textquotedblleft{}capable\textquotedblright{}, and Old High German \emph{gigat}, \textquotedblleft{}fitting\textquotedblright{}; the same root is in \emph{\ru{годиться}}, \textquotedblleft{}to be fit\textquotedblright{}.
\end{cousinweb}

\begin{grammarlens}[title={What you've built}]
One more word for this chapter.
\end{grammarlens}

\subsection*{Your turn}
\begin{itemize}
\raggedright
  \item \emph{Say it:} \emph{god}
  \item \emph{Say it:} \emph{god}, once more
  \item \emph{Say it:} say \emph{god}
  \item \emph{Recall:} say the Russian for a week, then the Russian for a month, then say \emph{god} again
\end{itemize}

\begin{tcolorbox}[breakable,colback=teal!4,colframe=teal!35!black,title={Before you move on}]
What does \ru{год} mean? (\textquotedblleft{}A year\textquotedblright{}.) Say it once more.
\end{tcolorbox}
\section[\texorpdfstring{\ru{вечер} (vécher)}{вечер (vécher)}]{\ru{вечер} (vécher) — the evening}

\label{lesson:RU-C68-evening}



\begin{quote}
Before the new one: say the Russian for a month, then the Russian for a year.
\end{quote}

\subsection*{You'll want to know: \ru{вечер}}
\textbf{\ru{вечер}} (\emph{vécher}) — \textquotedblleft{}the evening\textquotedblright{}.

\begin{cousinweb}
\textbf{\ru{вечер}} $\leftarrow$ Proto-Slavic \emph{*večerŭ} — inherited, a cousin of Lithuanian \emph{vãkaras}, \textquotedblleft{}evening\textquotedblright{} (Vasmer). Vasmer thinks Greek \emph{hespéra} and Latin \emph{vespera}, \textquotedblleft{}evening\textquotedblright{}, probably belong here too, with a changed first sound.
\end{cousinweb}

\begin{grammarlens}[title={What you've built}]
One more word for this chapter.
\end{grammarlens}

\subsection*{Your turn}
\begin{itemize}
\raggedright
  \item \emph{Say it:} \emph{vécher}
  \item \emph{Say it:} \emph{vécher}, once more
  \item \emph{Say it:} say \emph{vécher}
  \item \emph{Recall:} say the Russian for a month, then the Russian for a year, then say \emph{vécher} again
\end{itemize}

\begin{tcolorbox}[breakable,colback=teal!4,colframe=teal!35!black,title={Before you move on}]
What does \ru{вечер} mean? (\textquotedblleft{}The evening\textquotedblright{}.) Say it once more.
\end{tcolorbox}
